\documentclass[10pt,conference]{IEEEtran}

\usepackage{booktabs}
\usepackage{cite}
\usepackage[hidelinks]{hyperref}
\hypersetup{
  pdfauthor={Anonymous Authors},
  pdftitle={When Does Project Memory Help Coding Agents? A Fail-Closed Evaluation Protocol}
}

\newcommand{\bench}{\textsc{FreshMem}}

\title{When Does Project Memory Help Coding Agents?\\
  A Fail-Closed Evaluation Protocol}
\author{Anonymous Authors}

\begin{document}
\maketitle

\begin{abstract}
Project memory is increasingly added to coding agents, but a capable fresh
agent can often recover the needed state from the current checkout and tests.
Memory can therefore help, do nothing, or promote obsolete guidance. We present
\bench{}, a fail-closed evaluation protocol for locating that boundary. A transfer case
holds the evolved repository, model route, ordinary coding tools, and budget
fixed while varying only predecessor evidence: none, an evidence-capped raw
record, or native progressive-disclosure memory. The design treats an adapter
failure, route mismatch, or unequal tool surface as invalid rather than as a
task failure. An archived public pilot over twelve fully public fixtures is
inconclusive: one Qwen cohort reports a $+2.8$ percentage-point target-memory
difference versus no memory (95\% interval $[-8.3,16.7]$), while a DeepSeek
cohort reaches ceiling success for both conditions and shows added context cost.
A separate sealed-local runner has been exercised end-to-end, but is excluded
from efficacy claims. These observations motivate an evaluation protocol rather
than a claim that more memory is universally better.
\end{abstract}

\begin{IEEEkeywords}
agent memory, coding agents, evaluation, software evolution, reproducibility
\end{IEEEkeywords}

\section{Motivation}

Long-term-memory systems are usually justified by a simple intuition: an agent
that retains useful prior information should perform better later. This is often
reasonable, and existing work evaluates long-horizon retrieval and updates in
conversational settings~\cite{wu2024longmemeval}, practical long-term memory
architectures~\cite{chhikara2025mem0}, and incremental agent-memory
competencies~\cite{hu2025memoryagentbench}. Workflow-memory work further shows
that reusable routines can matter across agent runs~\cite{wang2024agentworkflowmemory}.

Coding agents create a different control problem. A later agent can inspect the
current repository, search symbols, edit files, and run focused tests. For a
source-sufficient task, predecessor context may merely consume tokens. For a
task governed by a prior design decision that source no longer explains, compact
memory may reduce ambiguity. After a repository change, a previously correct
claim may also become stale. The central question is therefore conditional:
\emph{under which task, freshness, and model-capability conditions does project
memory improve or harm a later coding task?}

This is not the first study of dynamic or sequential agent environments.
MemoryArena studies interdependent multi-session tasks~\cite{he2026memoryarena},
EvoArena studies evolving environments~\cite{xu2026evoarena}, and SWE-EVO and
ChainSWE study long-horizon software evolution~\cite{le2025sweevo,jin2026chainswe}.
\bench{} narrows the intervention: a fresh agent receives the same current
code and ordinary tools in every condition, while only bounded predecessor
evidence changes. Its contributions are:

\begin{enumerate}
  \item a fresh-agent transfer design that does not weaken the no-memory control;
  \item a fail-closed execution contract that separates task failures from
  broken routes, adapters, and unequal capabilities; and
  \item an archived public diagnostic whose mixed and ceiling outcomes motivate
  a boundary-focused study rather than a universal memory claim.
\end{enumerate}

\section{Fresh-Agent Boundary Design}

Each case has a precursor session, an intervening repository evolution, and a
transfer task for a fresh agent. The transfer checkout, task, writable paths,
ordinary inspection tools, verification command, timeout, and requested model
route are matched. The intervention is predecessor evidence, not access to
current source.

\begin{table}[t]
\caption{Matched transfer conditions. Every condition retains normal access to
the same current checkout and trusted verification.}\label{tab:conditions}
\centering
\small
\begin{tabular}{p{0.24\columnwidth}p{0.66\columnwidth}}
\toprule
Condition & Predecessor evidence before or during transfer \\
\midrule
No memory & None; current source and tests only \\
Raw record & One explicit read of a fixed predecessor record \\
Native memory & One bounded project brief and one cited detail expansion \\
\bottomrule
\end{tabular}
\end{table}

The memory-providing conditions share a case-defined evidence-size cap and
require an explicit tool call. The native condition models progressive disclosure
rather than injecting a full transcript: the agent first sees a bounded project
brief and can expand one cited memory if it is relevant. Provider-reported input
tokens remain a separate outcome, since routes need not tokenize identical text
identically. Native product behavior must not be simulated by adding
evaluation-only switches to the product interface.

The protocol separates two profiles rather than conflating their conclusions.
The native-product profile retains each condition's real tool surface and thus
measures total integration cost, including tool-schema overhead; the
canonical-information profile gives every condition the same neutral
predecessor-index and predecessor-detail tools; only the returned evidence is
none, a raw record, or the native backend. The latter isolates an information
effect, while the former retains product UX\@. Neither is a substitute for the
other.

Cases are classified before execution as source-sufficient controls,
durable-decision dependencies, or stale-conflict tasks. Model capability is a
moderator rather than a vendor label: each cohort freezes one provider-reported
route, and analysis asks whether effects differ across predeclared model tiers
and case classes. A ceiling result is informative; it can mean that the current
repository already supplies enough evidence for that setting.

\subsection{Fail-Closed Execution}

Patch correctness is the primary outcome and requires a deterministic verifier
after the agent exits. Token use, time, cost, tool calls, and stale actions are
reported beside correctness, not folded into it. A row is invalid when a
condition loses ordinary current-code capability, the workspace differs, the
model route is missing or mixed, the oracle changes, or the memory adapter
fails. Assigning such a failure a zero would manufacture a memory effect.

The exploratory runner gives the model only list, read, bounded write,
exact-text replacement, and trusted verification tools. It has no shell,
network, process, host-filesystem, or
oracle-file tool. The controller retains the private verifier outside the
agent-visible checkout and writes only hashes and sanitized events to an
external receipt. This is a useful local isolation boundary, not a claim of an
independent confirmatory sandbox.

\section{Initial Evidence and Scope}

Table~\ref{tab:pilot} reports an archived public pilot from an earlier target
implementation. The fixtures and checks were fully public, so the table cannot
establish that tasks were unseen in training or that a private-oracle effect
exists. The older Qwen cohort also used two non-target local adapters; they are
reported for completeness but were not configured under the new protocol. The
new raw-record control was not part of that pilot, so the table supports no
native-versus-raw claim. The target-versus-no-memory interval spans meaningful
benefit and harm, while the DeepSeek cohort has no correctness headroom.

\begin{table}[t]
\caption{Archived public pilot. Repetitions are collapsed within each of 12
fixture clusters; all quantities are descriptive.}\label{tab:pilot}
\centering
\small
\begin{tabular}{@{}lp{0.64\columnwidth}@{}}
\toprule
Cohort & Observation \\
\midrule
Qwen (144 rows) & No memory 34/36; local baseline one 34/36; local baseline two 35/36; target canonical 35/36. The predeclared target-versus-none contrast was $+2.8$ points (95\% interval $[-8.3,16.7]$), sign-flip $p=1.0$. \\
DeepSeek (72 rows) & No memory and target canonical: 36/36 each; target canonical added about 2,206 input tokens per cluster. \\
\bottomrule
\end{tabular}
\end{table}

The current sealed-local runner was additionally exercised on three synthetic
engineering cases under all three native-product conditions. All nine executions
completed, which is expected for small engineering cases and is not entered into
Table~\ref{tab:pilot} or used as an effectiveness result. Its purpose was
operational: it confirmed formation, native brief-to-detail retrieval, bounded
tools, route receipts, private verification, and sanitized artifact output can
work together without exposing the oracle to the model.

Several threats remain. Public coding tasks and fixes can enter training data,
as recent SWE-bench analysis emphasizes~\cite{openai2026swebench}. Provider
telemetry cannot prove the absence of opaque provider-side state. The public
pilot has only twelve fixture clusters and two routes, and the sealed-local
smoke is synthetic. Consequently, this paper makes no claim of general
effectiveness, cross-agent transfer, or superiority over any memory system.

\section{Future Plans}

The next stage is a preregistered confirmatory study. Each held-out case will
start from a permissively licensed source or a newly authored transition, then
bind a precursor record, private post-snapshot task, reference repair, and
behavioral oracle before any model outcome is read. Independent reviewers will
admit the dependency classification and reject cases whose current source makes
the task trivial. The worker will receive the evolved checkout and bounded
condition input; a separate controller will own the oracle.

Before collecting results, the study will freeze cases, model routes, budgets,
primary comparisons, and the required number of independent case clusters.
Each memory-providing condition will have the same predeclared evidence-size
cap; provider-reported input tokens will remain an outcome, not an assumed
match. A case-by-route-by-repetition is the paired analysis unit, and every row
will retain its completion, invalid, timeout, and cost status. A full paper can
then test whether the anticipated boundary persists across source-sufficient,
durable-decision, and stale-conflict cases and across frozen model capability
tiers. The resulting conclusion may be positive, negative, or mixed; the
design is intended to make each outcome interpretable.

\bibliographystyle{IEEEtran}
\bibliography{references}

\end{document}
